\documentclass[11pt]{article}

\usepackage[letterpaper,margin=1in]{geometry}
\usepackage[T1]{fontenc}
\usepackage[utf8]{inputenc}
\usepackage{lmodern}
\usepackage{microtype}
\usepackage{amsmath,amssymb,mathtools,bm}
\usepackage{slashed}
\usepackage{enumitem}
\usepackage{xcolor}
\usepackage{booktabs}
\usepackage{array}
\usepackage{hyperref}
\usepackage[most]{tcolorbox}
\usepackage{fancyhdr}
\usepackage{titlesec}
\usepackage{parskip}
\usepackage{needspace}

\definecolor{qftblue}{HTML}{245A8D}
\definecolor{qftgreen}{HTML}{2F6B4F}
\definecolor{qftred}{HTML}{A23B3B}
\definecolor{qftgray}{HTML}{F4F6F8}

\hypersetup{
  colorlinks=true,
  linkcolor=qftblue,
  urlcolor=qftblue,
  citecolor=qftgreen,
  pdftitle={PHYS 662 Lecture 16: Symmetries, Causality, and Propagators},
  pdfauthor={Nima Lashkari}
}

\pagestyle{fancy}
\fancyhf{}
\lhead{PHYS 662: Quantum Field Theory I}
\rhead{Lecture 16}
\cfoot{\thepage}
\setlength{\headheight}{14pt}

\titleformat{\section}{\Large\bfseries\color{qftblue}}{\thesection.}{0.6em}{}
\titleformat{\subsection}{\large\bfseries\color{qftgreen}}{\thesubsection}{0.6em}{}
\titleformat{\subsubsection}{\normalsize\bfseries\color{qftred}}{\thesubsubsection}{0.6em}{}

\setlist[itemize]{leftmargin=2em,itemsep=0.25em,topsep=0.4em}
\setlist[enumerate]{leftmargin=2em,itemsep=0.25em,topsep=0.4em}

\newcommand{\dd}{\mathrm{d}}
\newcommand{\R}{\mathbb{R}}
\newcommand{\Lag}{\mathcal{L}}
\newcommand{\cM}{\mathcal{M}}
\newcommand{\SO}{\mathrm{SO}}
\newcommand{\OO}{\mathrm{O}}

\newtcolorbox{keyidea}[1][]{
  enhanced, breakable, colback=qftgray, colframe=qftblue,
  boxrule=0.8pt, arc=1.5mm, title={#1}, fonttitle=\bfseries,
  coltitle=white
}
\newtcolorbox{remarkbox}[1][]{
  enhanced, breakable, colback=white, colframe=qftgreen,
  boxrule=0.7pt, arc=1.5mm, title={#1}, fonttitle=\bfseries,
  coltitle=white
}
\newtcolorbox{warningbox}[1][]{
  enhanced, breakable, colback=white, colframe=qftred,
  boxrule=0.7pt, arc=1.5mm, title={#1}, fonttitle=\bfseries,
  coltitle=white
}


\newcommand{\Th}{\Theta}
\newcommand{\Eul}{\mathcal E}
\newcommand{\Lie}{\mathscr L}

\newcommand{\C}{\mathbb C}
\newcommand{\Z}{\mathbb Z}
\newcommand{\SU}{\mathrm{SU}}
\newcommand{\UU}{\mathrm{U}}
\newcommand{\GL}{\mathrm{GL}}
\newcommand{\SL}{\mathrm{SL}}
\newcommand{\Sp}{\mathrm{Sp}}
\newcommand{\tr}{\operatorname{tr}}
\newcommand{\Ad}{\operatorname{Ad}}
\newcommand{\ad}{\operatorname{ad}}
\newcommand{\lieg}{\mathfrak g}
\newcommand{\lieh}{\mathfrak h}
\newcommand{\su}{\mathfrak{su}}
\newcommand{\so}{\mathfrak{so}}
\newcommand{\slie}{\mathfrak{sl}}


\newcommand{\cH}{\mathcal H}
\newcommand{\Spin}{\operatorname{Spin}}
\newcommand{\id}{\mathrm{id}}
\newcommand{\Lor}{\SO^+(1,3)}
\newcommand{\eps}{\varepsilon}
\newcommand{\one}{\mathbf 1}

\newcommand{\smu}{s}
\newcommand{\tmu}{\widetilde s}
\newcommand{\Eps}{E}
\newcommand{\pslash}{\slashed p}
\newcommand{\ISO}{\mathrm{ISO}}
\newcommand{\pL}{P_L}
\newcommand{\pR}{P_R}
\newcommand{\hc}{\mathrm{h.c.}}
\newcommand{\diag}{\operatorname{diag}}
\newcommand{\vol}{\mathrm{vol}}
\newcommand{\supp}{\operatorname{supp}}

\newcommand{\ddPi}{\dd\Pi}
\newcommand{\vac}{\lvert0\rangle}
\newcommand{\vbra}{\langle0\rvert}
\newcommand{\ii}{\mathrm i}
\newcommand{\Dc}{\mathcal C}
\newcommand{\Dp}{\mathcal P}
\newcommand{\Dt}{\mathcal T}
\begin{document}
\pagenumbering{gobble}
\begin{titlepage}
\centering
\vspace*{1.2in}
{\scshape\Large PHYS 662 --- Quantum Field Theory I\par}
\vspace{0.35in}
{\Huge\bfseries\color{qftblue} Lecture 16\par}
\vspace{0.18in}
{\LARGE\bfseries Symmetries, Causality, and Propagators\par}
\vspace{0.16in}
{\large\color{qftgreen} Discrete Transformations, Correlations, and Response\par}
\vspace{0.6in}
{\large Nima Lashkari\par}
{\large Fall 2026\par}
\vfill
\begin{minipage}{0.86\textwidth}\small
The quantum fields of Lecture 15 now become a laboratory for discrete symmetries and locality. We separate vacuum correlations from signals, and derive Green functions from their initial data and boundary conditions.

\medskip
Course page: \url{https://nimalashkari.com/quantum-field-theory-2026-part-i/}
\end{minipage}
\vspace{0.5in}
\end{titlepage}
\pagenumbering{roman}
\setcounter{tocdepth}{1}\begingroup\small\tableofcontents\endgroup
\newpage
\pagenumbering{arabic}

\section*{Roadmap}
\addcontentsline{toc}{section}{Roadmap}
Lecture 15 constructed the scalar and Dirac Fock spaces. We first determine
how charge conjugation, parity, and time reversal act on those states.
We then distinguish vacuum correlations from causal response, derive
retarded and advanced Green functions, and construct scalar and Dirac
Feynman propagators. Equal-time canonical data fix every source normalization.

\begin{keyidea}[Conventions and prerequisites]
Use \(\eta=\diag(1,-1,-1,-1)\), \(p\cdot x=E_pt-\bm p\cdot\bm x\),
and \(\dd\Pi_p=\dd^3p/[(2\pi)^3 2E_p]\). The spinor basis is
\[
\gamma^0=\begin{pmatrix}0&I\\I&0\end{pmatrix},\qquad
\gamma^i=\begin{pmatrix}0&-\sigma_i\\\sigma_i&0\end{pmatrix},\qquad
\bar\Psi=\Psi^\dagger\gamma^0.
\]
Lecture 15 fixes the mode algebra and the positive-energy particle and
antiparticle vacuum. The spin sums are those of Lecture 14:
\(\sum_su_s\bar u_s=\slashed p+m\),
\(\sum_sv_s\bar v_s=\slashed p-m\).
\end{keyidea}

\Needspace{7\baselineskip}
\section{Charge conjugation and parity on Fock space}
\subsection{The charge-conjugation matrix and its quantum implementer}
Keep separate the spinor matrix \(C\) and the unitary Hilbert-space
operator \(\Dc\). In the gamma basis recalled above, choose
\begin{equation}
 E=i\sigma_2=\begin{pmatrix}0&1\\-1&0\end{pmatrix},\qquad
 C=\begin{pmatrix}E&0\\0&-E\end{pmatrix}=i\gamma^0\gamma^2,
 \qquad C(\gamma^\mu)^T C^{-1}=-\gamma^\mu.
 \label{f-eq:Cmatrix}
\end{equation}
Then \(C^T=-C\), \(C^{-1}=-C\). Its overall phase is conventional.
The field transformation is
\begin{equation}
 \Dc\Psi(x)\Dc^{-1}=\Psi^c(x)=C\bar\Psi(x)^T
   =B_C(\Psi^\dagger(x))^T,\qquad B_C=C\gamma^0=-i\gamma^2=\begin{pmatrix}0&E\\-E&0\end{pmatrix}.
 \label{f-eq:Cfield}
\end{equation}
The transpose is on spinor indices. On a classical coefficient column
this is an antilinear map, but \(\Dc\) is unitary on quantum states:
\(\Dc i\Dc^{-1}=i\). Componentwise adjoints in \eqref{f-eq:Cfield} must not
be confused with antiunitarity of the implementing operator.

To verify the mode action, choose spinor phases so that
\begin{equation}
 v_s(p)=C\bar u_s(p)^T,\qquad u_s(p)=C\bar v_s(p)^T.
 \label{f-eq:Cspinors}
\end{equation}
For the square-root spinors of Lecture 14 this means
\(\eta_s=E\xi_s^*\). Substitution into the Dirac mode expansion of Lecture 15 exchanges its two
mode coefficients, giving
\begin{equation}
 \Dc b_s\Dc^{-1}=d_s,\qquad \Dc d_s\Dc^{-1}=b_s,\qquad
 \Dc Q\Dc^{-1}=-Q,\qquad \Dc H_N\Dc^{-1}=H_N.
 \label{f-eq:Cmodes}
\end{equation}
Other spin bases may insert unitary spin-label matrices and phases;
the exchange of charge sectors is invariant information.
Reordering the odd fields gives
\(\bar\Psi^c\Psi^c=\bar\Psi\Psi\) and
\(:\bar\Psi^c\gamma^\mu\Psi^c:=-:\bar\Psi\gamma^\mu\Psi:\).
For quantum composite operators these equalities use a consistent
normal-ordering or renormalization prescription.

\subsection{Parity and relative intrinsic parity}
Let \(x_P=(t,-\bm x)\). A parity symmetry of the free massive theory can act as
\begin{equation}
 \Dp\Psi(x)\Dp^{-1}=\eta_P\gamma^0\Psi(x_P),\qquad |\eta_P|=1.
 \label{f-eq:Pfield}
\end{equation}
Recall that the current is $j^\mu=\bar{\Psi}\gamma^\mu\Psi$.
The conditions \((\gamma^0)^\dagger\gamma^0=1\) and
\(\gamma^0\gamma^i\gamma^0=-\gamma^i\) implies that the probability/charge density \(j^0\) is even under parity whereas the spatial current \(j^i\) is odd,
at the reflected point, as expected from parity transform. They also establish invariance of the action.
Taking \(\eta_P=1\), the canonical-spin wavefunctions satisfy
\[
 \gamma^0u_s(p)=u_s(p_-),\qquad \gamma^0v_s(p)=-v_s(p_-),
 \quad p_-=(E_{\bm p},-\bm p).
\]
Comparison of modes gives
\begin{equation}
 \Dp b_s(\bm p)\Dp^{-1}=b_s(-\bm p),\qquad
 \Dp d_s(\bm p)\Dp^{-1}=-d_s(-\bm p).
 \label{f-eq:Pmodes}
\end{equation}
With the usual choice that parity squares to one on the field,
\(\eta_P=\pm1\), particle and antiparticle have opposite intrinsic
parities. Composing parity with an internal charge rotation changes
the definition of the symmetry and its phase assignments; that choice
must be specified separately. Parity invariance of the free Dirac action
does not guarantee parity invariance of arbitrary interactions.

\Needspace{7\baselineskip}
\section{Time reversal and CPT}
\subsection{Why time reversal is antiunitary}
A symmetry reversing the direction of time evolution must satisfy
\begin{equation}
 \Dt e^{-iHt}\Dt^{-1}=e^{+iHt}.
 \label{f-eq:Tevolution}
\end{equation}
If \(\Dt\) were unitary, this would require \(\Dt H\Dt^{-1}=-H\), which
is incompatible with a nontrivial spectrum bounded below and unbounded
above. Instead time reversal is antiunitary:
\begin{equation}
 \Dt i\Dt^{-1}=-i,\qquad \Dt H\Dt^{-1}=H,\qquad
 \Dt\bm P\Dt^{-1}=-\bm P.
 \label{f-eq:Tcharges}
\end{equation}
Energy is preserved. The sign of \(i\), not a sign reversal of physical
energy, reverses the exponential. Thus the spacetime reflection matrix
\(\diag(-1,1,1,1)\) cannot be applied to \(P^\mu\) by the same unitary
conjugation rule used for the connected Lorentz group.

\subsection{The spinor transformation and its phases}
For \(x_T=(-t,\bm x)\), one convenient choice in our basis is
\begin{equation}
 \Dt\Psi(x)\Dt^{-1}=B_T\Psi(x_T),\qquad B_T=i\gamma^1\gamma^3= \begin{pmatrix}i E&0\\0&i E\end{pmatrix}= \begin{pmatrix}-\sigma_2&0\\0&-\sigma_2\end{pmatrix}.
 \label{f-eq:Tfield}
\end{equation}
The useful matrix tests are
\begin{equation}
 B_T(\gamma^0)^*B_T^{-1}=\gamma^0,\qquad
 B_T(\gamma^i)^*B_T^{-1}=-\gamma^i,\qquad B_TB_T^*=-1.
 \label{f-eq:Tmatrix}
\end{equation}
They express the complex conjugation of numerical coefficients caused
by antiunitarity. On a classical solution the corresponding transformed
column is \(B_T\psi^*(-t,\bm x)\); on the field operator the conjugation
of coefficients is already performed by \(\Dt\).

For a canonical spin basis, the phases of states can be chosen so that
\[
 \Dt|\bm p,+\tfrac12\rangle=|-\bm p,-\tfrac12\rangle,\qquad
 \Dt|\bm p,-\tfrac12\rangle=-|-\bm p,+\tfrac12\rangle.
\]
The relative minus sign makes \(\Dt^2=-1\) on the one-fermion sector.
A rule that merely reverses momentum while leaving both spin labels
unchanged would miss spin reversal and this square. The antiparticle
sector has the same spin-reversal structure.

\subsection{Combining the operations}
With the specific phase choices above, \(\Theta=\Dc\Dp\Dt\) is antiunitary
and the successive field maps give
\begin{equation}
 \Theta\Psi(x)\Theta^{-1}
   =B_T\gamma^0B_C(\Psi^\dagger(-x))^T
   =-i\gamma_* (\Psi^\dagger(-x))^T,
 \quad \gamma_*= \begin{pmatrix}-\mathbb{I}&0\\0&\mathbb{I}\end{pmatrix} =-i\gamma^0\gamma^1\gamma^2\gamma^3.
 \label{f-eq:CPT}
\end{equation}
The phase and the sign used to name \(\gamma_*\) are conventional.
The physical operation exchanges particle and antiparticle, reverses
spin, and leaves the momentum label unchanged because parity and time
reversal each reverse it.

\begin{remarkbox}[The scope of the CPT statement]
The free theory here has \(C\), \(P\), and \(T\) separately. More generally,
local relativistic quantum field theories satisfying the usual
Lorentz-covariance, positive-energy, and positive-Hilbert-space assumptions
obey the CPT theorem, even when the separate operations fail to be
symmetries. Locality by itself, without the other hypotheses, is not
a sufficient statement of the theorem. The explicit free-field matrices
illustrate CPT; they do not constitute its general proof.
\end{remarkbox}



\Needspace{7\baselineskip}
\section{Locality and the scalar commutator}
\begin{keyidea}[Fix the metric and every factor of \(i\)]
Use \(\eta=\diag(1,-1,-1,-1)\), \(z=x-y=(t,\bm r)\), and
\(p\cdot z=E_{\bm p}t-\bm p\cdot\bm r\). Thus spacelike means
\(z^2=t^2-r^2<0\). For a real scalar field define
\[
 D^+(z)=\vbra\phi(x)\phi(y)\vac,\qquad
 \mathcal K(z)=[\phi(x),\phi(y)]=D^+(z)-D^+(-z).
\]
The last equality identifies a c-number distribution multiplying the
identity operator. We suppress that identity. Our \(\mathcal K\) is purely
imaginary for real arguments; a real Pauli--Jordan distribution can be
defined by \(\Delta=\mathcal K/i\). All Green functions called \(G\) below
are normalized to \((\Box+m^2)G=\delta^{(4)}\).
\end{keyidea}
\subsection{From equal time to spacelike separation}
Canonical quantization gives \([\phi(t,\bm x),\phi(t,\bm y)]=0\).
Any two spacelike-separated events can be brought to equal time by a
proper orthochronous Lorentz transformation. Because the field transforms
covariantly, the operator commutator must vanish for that pair in every
inertial frame:
\begin{equation}
 [\phi(x),\phi(y)]=0\qquad\text{if}\qquad (x-y)^2<0.
 \label{k-eq:locality}
\end{equation}
This is microcausality. More generally, physical local observables
commute when their supports are spacelike separated. Fermionic odd
fields obey the corresponding graded relation, discussed below.
Derivatives of a commutator vanish in the open spacelike region as well;
there may still be distributions supported on the light cone.

The mode expansion supplies a direct verification:
\begin{equation}
 \mathcal K(z)=\int\frac{\dd^3p}{(2\pi)^3 2E_{\bm p}}
 (e^{-ipz}-e^{ipz}).
 \label{k-eq:commutator}
\end{equation}
In a frame with \(t=0\), the two integrals are equal after
\(\bm p\mapsto-\bm p\). The invariant mass-shell measure extends this
cancellation to every spacelike \(z\). Equivalently, a connected Lorentz
transformation followed by a spatial rotation takes any spacelike
\(z\) to \(-z\); this is not possible for a future-directed timelike
vector without changing time orientation.

\subsection{Initial data fix the normalization}
Each plane wave is on shell, so
\((\Box+m^2)\mathcal K=0\). At equal time,
\begin{equation}
 \mathcal K(0,\bm r)=0,\qquad
 \partial_t\mathcal K(0,\bm r)=-i\delta^{(3)}(\bm r).
 \label{k-eq:initialdata}
\end{equation}
The second sign follows either by differentiating \eqref{k-eq:commutator}
or from \([\dot\phi(0,\bm x),\phi(0,\bm y)]=-i\delta^{(3)}(\bm x-\bm y)\).
These initial data will determine the delta-function sources of the
retarded and advanced Green functions without an ambiguous contour sign.

For the complex scalar, we had the mode expansion
\begin{align}
	\Phi(x)=\int d\Pi_p(a(p)e^{-i p\cdot x}+b^\dagger(p) e^{i p\cdot x})
\end{align}
with $a(p)$ and $b(p)$ the annihilation modes for, respectively, the particle and the anti-particle.
Then, the commutator becomes
\begin{equation}
 [\Phi(x),\Phi^\dagger(y)]=\mathcal K(x-y),\qquad
 [\Phi(x),\Phi(y)]=0.
 \label{k-eq:complexcomm}
\end{equation}
The first term of \(\mathcal K\) comes from the particle commutator and the
second from the antiparticle commutator, whose reversed order gives its
minus sign. Keeping only the positive-frequency particle term would
leave a nonlocal kernel. Antiparticles are part of the local relativistic
field construction, even though their excitations have positive energy.

\Needspace{7\baselineskip}
\section{Vacuum correlations do not vanish outside the light cone}
\subsection{The Wightman function}
Only an annihilator followed by a creator has a nonzero vacuum expectation,
so
\begin{equation}
 D^+(z)=\int\dd\Pi_p\,e^{-ipz},\qquad
 D^+(-z)=D^+(z)^*.
 \label{k-eq:Wightman}
\end{equation}
This is called a Wightman function.
For a complex scalar this is
\(\vbra\Phi(x)\Phi^\dagger(y)\vac\); the charge-violating correlators
\(\vbra\Phi(x)\Phi(y)\vac\) and its adjoint vanish in the charge-invariant
free vacuum. 

The integral is an oscillatory distribution. The term $e^{-i p z}$ oscillates endlessly at large energies and the integral does not strictly converge. To fix this, since $E_p>0$ we introduce a real damping term $e^{-E_p}$, by shifting time slightly into the complex plane by letting $t\to t-i \epsilon$ with $\epsilon>0$. The answer we find following this prescription and sending $\epsilon\to 0^+$ is called the  positive-frequency ($E_p>0$)
boundary value $(\epsilon\to 0^+)$. Lorentz invariance means invariance under the
proper orthochronous group, preserving the future mass shell. At timelike
separation the sign of \(t\) remains relevant because restricted Lorentz transformations cannot change the sign of $t$. Therefore, writing a function of
\(z^2\) without a boundary prescription loses information. We need an $i\epsilon$-prescription.

\subsection{Spacelike evaluation and the Compton scale}
For spacelike separation choose \(t=0\) and let \(r=\sqrt{-z^2}>0\).
Angular integration gives
\begin{equation}
 D^+(0,\bm r)=\frac{1}{4\pi^2r}\int_0^\infty\dd p\,
                 \frac{p\sin(pr)}{\sqrt{p^2+m^2}}
              =\frac{m}{4\pi^2r}K_1(mr).
 \label{k-eq:Bessel}
\end{equation}
Here \(K_1\) is the modified Bessel function of the second kind (do not confused it with a
Lorentz boost generator). One way to verify the coefficient is to use
\[
 \frac{1}{2\sqrt{p^2+m^2}}
 =\frac{1}{2\sqrt\pi}\int_0^\infty\dd\alpha\,\alpha^{-1/2}
                         e^{-\alpha(p^2+m^2)}.
\]
The Gaussian momentum integral yields (work it out)
\begin{align}
	\int_0^\infty dp \:p\sin(p r)e^{-\alpha p^2}&=\frac{1}{2}\Im\int_{-\infty}^\infty dp\: p e^{-\alpha p^2+i p r}\nonumber\\
	&=\frac{r \sqrt{\pi}}{4\alpha^{3/2}}e^{-r^2/(4\alpha)}
\end{align}
which plugged back into \eqref{k-eq:Bessel} gives an integral representation of the Wightman function:
\begin{equation}
 D^+(0,\bm r)=\frac{1}{16\pi^2}\int_0^\infty\dd\alpha\,
 \alpha^{-2}e^{-m^2\alpha-r^2/(4\alpha)},
 \label{k-eq:Schwinger}
\end{equation}
This also makes positivity at nonzero spacelike distance manifest.

For \(mr\gg1\), \(K_1(mr)\sim\sqrt{\pi/(2mr)}e^{-mr}\), so
\begin{equation}
 D^+(0,\bm r)\sim\sqrt{\frac{m}{32\pi^3r^3}}\,e^{-mr}.
 \label{k-eq:spacelikeasymptotic}
\end{equation}
The scale \(m^{-1}\) is the Compton wavelength. In the massless limit at
fixed nonzero \(r\), \(K_1(mr)\sim1/(mr)\), giving
\(D^+=1/(4\pi^2r^2)\). The correlation then falls as a power rather than
exponentially. These limits are different from the singular coincidence
limit \(r\to0\).

\subsection{Timelike behavior}
For timelike separation choose \(\bm r=0\). In three spatial dimensions,
\begin{equation}
 D^+(t,0)=\frac{1}{4\pi^2}\int_m^\infty\dd E\,
               \sqrt{E^2-m^2}\,e^{-iE(t-i0^+)}.
 \label{k-eq:timelike}
\end{equation}
For \(mt\gg1\), \(t>0\), the stationary region is \(|\bm p|\ll m\).
Using \(E_{\bm p}=m+\bm p^2/(2m)+\cdots\) in the mode integral gives
\begin{equation}
 D^+(t,0)\sim\frac{e^{-imt}}{2m}
                   \left(\frac{m}{2\pi i t}\right)^{3/2}.
 \label{k-eq:timelikeasymptotic}
\end{equation}
It oscillates with a \(t^{-3/2}\) envelope; it does not decay as
\(e^{-m|t|}\). For \(t<0\), it follows from \eqref{k-eq:Wightman} that we only to need to take the complex conjugate of \(t>0\). A compact formula
containing both regimes is
\begin{equation}
 D^+(t,\bm r)=\lim_{\epsilon\downarrow0}
 \frac{m}{4\pi^2\rho_\epsilon}K_1(m\rho_\epsilon),\qquad
 \rho_\epsilon=\sqrt{r^2-(t-i\epsilon)^2},
 \label{k-eq:analyticW}
\end{equation}
with the square root continued from positive real \(\rho\).

\Needspace{7\baselineskip}
\section{Locality, statistics, and signaling}
\subsection{Why a nonzero correlation is not a signal}
At spacelike separation, \(D^+(z)=D^+(-z)\) and both can be nonzero,
while their difference \(\mathcal K(z)\) vanishes. A correlation measures
properties of the state. A controllable influence depends on a commutator.
For example, let \(A\) be an observable in one region, \(B\) a Hermitian
observable in a spacelike-separated region, and \(U_B=e^{-i\lambda B}\).
Then
\[
 \langle U_B^\dagger A U_B\rangle
 =\langle A\rangle+i\lambda\langle[B,A]\rangle+O(\lambda^2).
\]
If \([A,B]=0\), the expectation is unchanged to every order. Vacuum
entanglement can give correlations between the regions without enabling
this local operation to send a message.

\subsection{The free scalar test of statistics}
Take a positive-norm scalar mode family and a trial field
\(F_\lambda(x)=\int\dd\Pi_p[a(\bm p)e^{-ipx}
+\lambda a^\dagger(\bm p)e^{ipx}]\).
At spacelike separation its bosonic commutator with its adjoint is
\((1-|\lambda|^2)D^+(z)\), which vanishes for \(|\lambda|=1\).
If the same modes instead obey CAR, its anticommutator is
\((1+|\lambda|^2)D^+(z)\), which cannot vanish for any \(\lambda\).
Even the \(\lambda=0\) positive-frequency field is nonlocal.
This calculation identifies the cancellation needed for a local free
scalar field. A general spin--statistics theorem requires further
hypotheses, including relativistic covariance, energy positivity, and
a positive Hilbert-space inner product.

For the Dirac field the local relation is an anticommutator:
\begin{equation}
 \{\Psi_\alpha(x),\bar\Psi_\beta(y)\}
   =(i\gamma^\mu\partial_\mu+m)_{\alpha\beta}\mathcal K(x-y),
 \qquad \{\Psi_\alpha(x),\Psi_\beta(y)\}=0.
 \label{k-eq:Diraclocal}
\end{equation}
Use the two spin sums from Lecture 15 to verify this: differentiating the
negative-frequency term reverses the sign of \(\slashed p\), while the
minus in \(\mathcal K\) reverses it again. Since \(\mathcal K\) vanishes in
the open spacelike region, so does the anticommutator. Even local
bilinears commute at spacelike separation: moving each of two fermions
past the two in the other bilinear gives four minus signs. Coincident
products still require their usual regulator or subtraction.

\Needspace{7\baselineskip}
\section{Retarded and advanced Green functions}
\subsection{Derivation from the equal-time jump}
Let \(L=\Box+m^2\). Define
\begin{equation}
 G_R(z)=i\theta(t)\mathcal K(z),\qquad
 G_A(z)=-i\theta(-t)\mathcal K(z).
 \label{k-eq:RA}
\end{equation}
The opposite signs are essential. For a homogeneous solution \(f\) with
\(f(0,\bm r)=0\), distributional differentiation gives
\[
 L[\theta(t)f]=\delta(t)\partial_t f(0,\bm r),\qquad
 L[\theta(-t)f]=-\delta(t)\partial_t f(0,\bm r).
\]
The apparent \(\delta' f\) and \(2\delta\dot f\) terms combine into one
jump term; treating them as ordinary functions would give a spurious
factor of two. Applying \eqref{k-eq:initialdata} gives
\begin{equation}
 LG_R=LG_A=\delta^{(4)}(z),\qquad G_R-G_A=i\mathcal K.
 \label{k-eq:RAequation}
\end{equation}
The retarded function has support in the future light cone; the advanced
function in the past light cone. The commutator itself has both supports
and is a homogeneous solution, not a retarded Green function.

In spatial Fourier space these facts are especially transparent:
\begin{equation}
 G_R(t,\bm p)=\theta(t)\frac{\sin(E_{\bm p}t)}{E_{\bm p}},\qquad
 G_A(t,\bm p)=-\theta(-t)\frac{\sin(E_{\bm p}t)}{E_{\bm p}}.
 \label{k-eq:RAoscillator}
\end{equation}
Both have a unit jump in their first derivative across \(t=0\).

\subsection{Response to a local source}
Add \(J\phi\) to the Lagrangian, so that the equation of motion becomes \(L\phi=J\). With no incoming change
in the initial data, the induced field is
\begin{equation}
 \delta\phi(x)=\int\dd^4y\,G_R(x-y)J(y).
 \label{k-eq:response}
\end{equation}
The response depends only on sources in the causal past of \(x\).
Linear response in the quantum theory gives the same result because
\(G_R=i\theta(x^0-y^0)\langle[\phi(x),\phi(y)]\rangle\).
This is the operational role of the commutator that a Wightman
correlation alone does not have.

In \(3+1\) dimensions a useful explicit check is
\begin{equation}
 G_R(t,r)=\theta(t)\left[
 \frac{\delta(t^2-r^2)}{2\pi}
 -\theta(t^2-r^2)\frac{mJ_1(m\sqrt{t^2-r^2})}{4\pi\sqrt{t^2-r^2}}
 \right].
 \label{k-eq:retardedexplicit}
\end{equation}
The first term propagates on the cone. A massive field also has an
interior tail. For \(m=0\) the tail disappears and
\(G_R=\delta(t-r)/(4\pi r)\) for \(r>0\). The light-cone delta term cannot
be recovered by evaluating the Bessel expression only away from the cone.

\Needspace{7\baselineskip}
\section{Pole prescriptions encode boundary conditions}
\subsection{Retarded and advanced poles}
Fourier transforming \(LG=\delta^{(4)}\) suggests
\(\widetilde G=-1/(p^2-m^2)\), but this is singular on the mass shell.
The boundary condition prescribes how to pass the poles:
\begin{align}
 G_R(z)&=\lim_{\epsilon\downarrow0}\int\frac{\dd^4p}{(2\pi)^4}
 \frac{-e^{-ipz}}{(p^0+i\epsilon)^2-E_{\bm p}^2},\notag\\
 G_A(z)&=\lim_{\epsilon\downarrow0}\int\frac{\dd^4p}{(2\pi)^4}
 \frac{-e^{-ipz}}{(p^0-i\epsilon)^2-E_{\bm p}^2}.
 \label{k-eq:RApoles}
\end{align}
For the retarded prescription both poles are below the real axis.
When \(t<0\), close in the upper half-plane, where
\(e^{-ip^0t}\) decays: there are no poles, so the answer is zero.
When \(t>0\), close clockwise below; the residues give
\(\sin(E_{\bm p}t)/E_{\bm p}\). Reversing the displacements gives the
advanced prescription. Contour orientation supplies a minus sign and
must be kept together with the numerator in \eqref{k-eq:RApoles}.

The equation alone cannot choose a prescription. The difference of two
Green functions solves the homogeneous equation; boundary or initial
conditions provide the missing information.

\subsection{Time ordering and the Feynman function}
The scalar time-ordered vacuum expectation is
\begin{equation}
 D_F(z)=\vbra T\phi(x)\phi(y)\vac
 =\theta(t)D^+(z)+\theta(-t)D^+(-z).
 \label{k-eq:DFtime}
\end{equation}
Its Fourier representation is
\begin{equation}
 D_F(z)=\lim_{\epsilon\downarrow0}\int\frac{\dd^4p}{(2\pi)^4}
        \frac{i\,e^{-ipz}}{p^2-m^2+i\epsilon}.
 \label{k-eq:DFfourier}
\end{equation}
The positive-energy pole lies just below and the negative-energy pole
just above the real axis. Closing below for \(t>0\) keeps the
positive-frequency mode; closing above for \(t<0\) keeps the
negative-frequency contribution. This gives \eqref{k-eq:DFtime}, including
its normalization.

Unlike \(\mathcal K\), time ordering has a derivative jump:
\begin{equation}
 LD_F=-i\delta^{(4)}(z),\qquad
 G_F\equiv iD_F,\qquad LG_F=\delta^{(4)}(z).
 \label{k-eq:DFsource}
\end{equation}
The name ``Feynman propagator'' is used for either \(D_F\), \(G_F\), or a
function differing by a factor of \(i\) in different texts. The defining
expectation value and source equation remove that ambiguity.
\begin{center}\small
\begin{tabular}{llll}
\toprule
Object & Definition & \(L\) acting on it & Support/property\\
\midrule
\(D^+\) & \(\langle\phi(x)\phi(y)\rangle\) & \(0\) & positive frequency\\
\(\mathcal K\) & \(D^+(z)-D^+(-z)\) & \(0\) & causal cone\\
\(G_R\) & \(i\theta(t)\mathcal K\) & \(\delta^{(4)}\) & future cone\\
\(G_A\) & \(-i\theta(-t)\mathcal K\) & \(\delta^{(4)}\) & past cone\\
\(D_F\) & \(\langle T\phi(x)\phi(y)\rangle\) & \(-i\delta^{(4)}\) & time ordering\\
\bottomrule
\end{tabular}
\end{center}
The Feynman function is nonzero at spacelike separation. This is
consistent with causality: it encodes vacuum time ordering, not the
retarded response to a controllable source. It is the two-point kernel
used in scattering and path-integral perturbation theory.

\Needspace{7\baselineskip}
\section{The Dirac Feynman propagator}

For a Dirac field $\Psi$, inserting the mode expansion into the Wightman functions we find
\begin{align}
&	\vbra\Psi_\alpha(x)\bar{\Psi}_\beta(y)\vac=\int d\Pi_p (\sum_s u_s(p)\bar{u}_s(p))_{\alpha\beta}e^{-i p(x-y)}=\int d\Pi_p(\slashed p+m)_{\alpha\beta} e^{-i p(x-y)}\nonumber\\
	&	\vbra\bar{\Psi}_\alpha(x)\Psi_\beta(y)\vac=\int d\Pi_p (\sum_s v_s(p)\bar{v}_s(p))_{\beta\alpha}e^{-i p(x-y)}=\int d\Pi_p(\slashed p-m)_{\beta\alpha} e^{-i p(x-y)}
\end{align}

Fermionic time ordering includes the exchange sign:
\begin{equation}
 S_F(x-y)_{\alpha\beta}=\vbra T\Psi_\alpha(x)\bar\Psi_\beta(y)\vac
 =\theta(t)\vbra\Psi_\alpha(x)\bar\Psi_\beta(y)\vac
 -\theta(-t)\vbra\bar\Psi_\beta(y)\Psi_\alpha(x)\vac.
 \label{k-eq:SFtime}
\end{equation}
Insert the mode expansion and the spin sums. For \(t>0\) the numerator is
\(\slashed p+m\); for \(t<0\) the exchange sign combines with the
antiparticle spin sum. The result is
\begin{equation}
 S_F(z)=(i\gamma^\mu\partial_\mu+m)D_F(z)
 =\int\frac{\dd^4p}{(2\pi)^4}
       \frac{i(\slashed p+m)e^{-ipz}}{p^2-m^2+i0^+}.
 \label{k-eq:SF}
\end{equation}
Since
\((i\slashed\partial-m)(i\slashed\partial+m)=-(\Box+m^2)\),
\begin{equation}
 (i\slashed\partial-m)S_F(z)=i\delta^{(4)}(z)I_4.
 \label{k-eq:SFsource}
\end{equation}
The scalar denominator fixes propagation and its boundary condition;
the numerator carries the spinor structure. This links the invariant
tensors of Lectures 12--13, the spin sums of Lecture 14, and the
quantization of Lecture 15 without repeating those derivations.

\section*{Summary}
\addcontentsline{toc}{section}{Summary}
Charge conjugation exchanges charge sectors; parity reverses momentum;
antiunitary time reversal preserves energy and reverses momentum and spin.
Their implementation must respect the mode phases and spinor basis.
Vacuum correlations can be nonzero at spacelike separation, while local
observables still commute there. Retarded and advanced Green functions
encode causal support and have opposite step-function prefactors.
Feynman time ordering chooses a different boundary condition. The Dirac
propagator follows by applying its numerator to the scalar Feynman function.
The canonical equal-time jump fixes all delta-function sources.

\newpage
\section*{Questions for review}
\addcontentsline{toc}{section}{Questions for review}
\begin{enumerate}
\item Verify the four matrix identities for \(C\) in \eqref{f-eq:Cmatrix} and
the spinor relation \eqref{f-eq:Cspinors}.
\item Derive the opposite intrinsic parities directly from the mode expansion.
\item Explain why time reversal preserves \(H\) but conjugates \(i\).
Check \(B_TB_T^*=-1\).
\item Distinguish a classical antilinear spinor map, a unitary charge
conjugation operator, and an antiunitary time reversal operator.
\item Prove microcausality both from equal-time covariance and from the
mode integral. Why must the separation be spacelike?
\item Derive \(\partial_t\mathcal K(0,\bm r)=-i\delta^{(3)}(\bm r)\).
\item Reproduce the coefficient of \(K_1(mr)\) using the Gaussian integral.
Compare large spacelike and large timelike separation.
\item Why can \(D^+(z)\neq0\) coexist with the impossibility of signaling?
\item Show why a positive-frequency scalar field by itself is nonlocal.
\item Verify the Dirac anticommutator and explain how even observables
can commute when the underlying fields anticommute.
\item Differentiate the step-function definitions of \(G_R,G_A\) and
track their relative sign.
\item Locate all poles in the retarded, advanced, and Feynman prescriptions.
Explain which half-plane is used for each sign of time.
\item Why is the commutator a homogeneous solution while a propagator
with a step function can have a delta-function source?
\item Verify the source equations for \(D_F\), \(G_F\), and \(S_F\).
\end{enumerate}

\section*{Sources and reading}
\addcontentsline{toc}{section}{Sources and reading}
Based on the discrete-symmetry sections of the supplied Dirac and
causality slides, the scalar locality/statistics example in the Dirac
slides, and the complete propagator discussion in the causality slides.
Compare D.~Tong, \emph{Quantum Field Theory}, Chapters 2 and 5,
\url{https://www.damtp.cam.ac.uk/user/tong/qft.htm}.
Bessel integral representations and asymptotics are collected in
\url{https://dlmf.nist.gov/10.32} and \url{https://dlmf.nist.gov/10.40}.
\end{document}
